%! Unit Launch: Study Design — Unit 1, Lesson 1.0 — Guided Notes
\documentclass[10pt]{article}
\usepackage{saar-article}
\usepackage{saar-boxes}
\newcommand{\TallMath}[1]{$\displaystyle #1\rule[-1.4em]{0pt}{3.2em}$}

\begin{document}

\pageheader{Unit 1, Lesson 1.0}{Guided Notes}
% NAMESTRIP: no \namedateperiod here. The student writes their name once, on the
% cover this component is stapled behind. See references/conventions.md.

\vspace{0.15in}

\begin{objectivebox}
\small
By the end of this lesson, I will be able to\dots
\begin{enumerate}[leftmargin=1.5em, itemsep=2pt, topsep=2pt]
  \item explain why a number is not data until I know \emph{who} was measured,
        \emph{what} was measured, and in \emph{what units};
  \item read a data table and name its \textbf{individuals} and its
        \textbf{variables};
  \item classify a variable as \textbf{categorical} or \textbf{quantitative},
        and justify my choice;
  \item recognize a variable written with digits that is really a label, and
        explain how I know;
  \item tell a \textbf{statistical question} from a question that is not
        statistical, and rewrite one so that it becomes statistical.
\end{enumerate}
\end{objectivebox}

\boxguard
\begin{vocabbox}
\small
Fill in each term as we build it together.
\termblanklong{Data}
\termblanklong{Individual}
\termblanklong{Variable}
\termblanklong{Categorical variable}
\termblanklong{Quantitative variable}
\termblanklong{Statistical question}
\end{vocabbox}

\boxguard[12]
\begin{hookbox}
\small
A poster in the Riverbend hallway says:

\begin{center}
\textit{``Four out of five students at Riverbend High School say the lunch line
is too slow.''}
\end{center}

\textbf{Do you believe it?} Write down \emph{one} question you would want
answered before you decided.
\writeline
\writeline
\end{hookbox}

\boxguard
\begin{notesbox}{1. Data Are Numbers With a Context}
\small
Write this number down: \qquad \textbf{\large 18}

\textbf{What does it tell you?} \blank{9cm}

Now here is the same number \emph{with its context:}

\begin{center}
\textit{``18 minutes --- the mean one-way commute for the six Riverbend students
we surveyed.''}
\end{center}

Three things had to be added before the number meant anything:

\begin{enumerate}[leftmargin=1.5em, itemsep=3pt, topsep=2pt]
  \item \textbf{Who} was measured? \blank{7.5cm}
  \item \textbf{What} was measured? \blank{7.5cm}
  \item In what \textbf{units}? \blank{7.5cm}
\end{enumerate}

\vspace{3pt}
\textbf{Big idea:} \blank{4.2cm} are numbers or labels \emph{together with the
context that gives them meaning}. A number with no context is not data --- it is
just a number.
\end{notesbox}

\boxguard[30]
\begin{notesbox}{2. Individuals and Variables}
\small
Six Riverbend students filled out the Student Life Survey.

\begin{center}
\renewcommand{\arraystretch}{1.2}
\begin{tabular}{@{} l c c l c c @{}}
\toprule
\textbf{Student} & \textbf{Grade} & \textbf{Commute (min)} & \textbf{Lunch} &
\textbf{Siblings} & \textbf{Club member} \\
\midrule
Alina   & 10 & 20 & Home       & 2 & Yes \\
Devon   & 11 &  8 & Cafeteria  & 0 & No  \\
Marisol &  9 & 35 & Home       & 3 & Yes \\
Theo    & 12 & 15 & Cafeteria  & 1 & Yes \\
Priya   & 10 &  5 & Off campus & 2 & No  \\
Jamal   & 11 & 25 & Cafeteria  & 4 & Yes \\
\bottomrule
\end{tabular}
\end{center}

\vspace{3pt}
Each \textbf{row} is one \blank{3.6cm} --- one person the data are about.

Each \textbf{column} is one \blank{3.6cm} --- one characteristic recorded for
every individual.

\vspace{3pt}
How many individuals are in this table? \blank{2cm}

How many variables are in this table? \blank{2cm}

\vspace{2pt}
\textit{Careful: the ``Student'' column names the individuals. It is a label for
the rows, not a variable we are studying.}
\end{notesbox}

\boxguard
\begin{notesbox}{3. The Two Types of Variables}
\small
Look at two columns side by side and ask one question:
\textbf{could I find an average of this column?}

\begin{center}
\renewcommand{\arraystretch}{1.3}
\begin{tabularx}{0.95\linewidth}{@{} p{3.4cm} X @{}}
\toprule
\textbf{Commute (min)} & $20, \; 8, \; 35, \; 15, \; 5, \; 25$ \\
Average it? & \blank{6cm} \\
\midrule
\textbf{Lunch} & Home, Cafeteria, Home, Cafeteria, Off campus, Cafeteria \\
Average it? & \blank{6cm} \\
\bottomrule
\end{tabularx}
\end{center}

\vspace{4pt}
That contrast \emph{is} the definition:

\begin{itemize}[leftmargin=1.3em, itemsep=3pt, topsep=3pt]
  \item A \textbf{quantitative variable} records a \blank{4.5cm} with units.\\
        Summarize it with a \blank{3cm}.
  \item A \textbf{categorical variable} records a \blank{4.5cm} or group name.\\
        Summarize it with \blank{3cm} and \blank{3cm}.
\end{itemize}

\vspace{3pt}
Classify each variable from the survey table:

\begin{center}
\renewcommand{\arraystretch}{1.45}
\begin{tabularx}{0.95\linewidth}{@{} p{3.2cm} p{3.4cm} X @{}}
\toprule
\textbf{Variable} & \textbf{Type} & \textbf{How I know} \\
\midrule
Commute (min) & \blank{2.8cm} & \blank{5cm} \\
Lunch         & \blank{2.8cm} & \blank{5cm} \\
Siblings      & \blank{2.8cm} & \blank{5cm} \\
Club member   & \blank{2.8cm} & \blank{5cm} \\
\bottomrule
\end{tabularx}
\end{center}
\end{notesbox}

\boxguard
\begin{notesbox}{4. Careful --- Numbers That Are Really Labels}
\small
We left one variable out on purpose: \textbf{Grade}, recorded as
$10, \; 11, \; 9, \; 12, \; 10, \; 11$.

It is made of digits. \textbf{Does that make it quantitative?} Let's test it by
actually computing the ``mean grade level.''

\begin{work}
  \bar{x} &= \dfrac{10+11+9+12+10+11}{6} \\
          &= \dfrac{63}{6} \\
          &= 10.5
\end{work}

\textbf{Who is in grade $10.5$?} \blank{8cm}

\vspace{2pt}
So Grade is a \blank{4cm} written with digits --- the type is
\blank{3.5cm}.

\vspace{4pt}
\textbf{The test to write down and keep:}
\begin{center}
\fbox{\parbox{0.88\linewidth}{\centering
Not ``does it look like a number?'' but\\[2pt]
\textbf{does arithmetic on this variable mean anything?}}}
\end{center}

\vspace{3pt}
Two more variables that fail the test even though they are all digits:

\hspace{1.2em} \blank{5cm} \hfill and \hfill \blank{5cm} \hspace{1.2em}
\end{notesbox}

\boxguard
\begin{notesbox}{5. Statistical Questions}
\small
\textbf{``How tall am I?''} --- measure once and you are done. One answer.

\textbf{``How tall are students at Riverbend?''} --- the answers
\blank{4cm} from student to student, so it takes data to answer.

\vspace{3pt}
A \textbf{statistical question} is one whose answers are expected to
\blank{3.5cm}.

\vspace{4pt}
Mark each question \textbf{S} (statistical) or \textbf{N} (not statistical):

\begin{enumerate}[leftmargin=1.5em, itemsep=4pt, topsep=3pt]
  \item How many students are enrolled at Riverbend? \hfill \blank{1.2cm}
  \item How long is the commute for students at Riverbend? \hfill \blank{1.2cm}
  \item What did Devon eat for lunch today? \hfill \blank{1.2cm}
  \item Where do Riverbend students eat lunch? \hfill \blank{1.2cm}
\end{enumerate}

\vspace{3pt}
\textit{Notice: a question can contain a number and still not be statistical.
What matters is whether different individuals give different answers.}

\vspace{3pt}
Rewrite question 3 so that it \emph{becomes} statistical:
\writeline
\end{notesbox}

\boxguard
\begin{notesbox}{6. Road Map --- Where Unit 1 Is Going}
\small
Every study in this course runs the same four stages, over and over. This is the
\textbf{data cycle}:

\begin{center}
\renewcommand{\arraystretch}{1.3}
\begin{tabularx}{0.95\linewidth}{@{} c p{4.2cm} X @{}}
\toprule
\textbf{Stage} & \textbf{What you do} & \textbf{What can go wrong} \\
\midrule
1 & Formulate the question   & The question is not statistical, or nobody says
                               who it is about. \\
2 & Collect or acquire data  & You ask the wrong people, or too few, or in a way
                               that pushes an answer. \\
3 & Organize and represent   & The display hides the pattern instead of showing
                               it. \\
4 & Analyze and communicate  & The conclusion claims more than the data support. \\
\bottomrule
\end{tabularx}
\end{center}

\vspace{3pt}
Go back to the lunch-line poster from the Hook. Which stage was your objection
attacking? \blank{2cm}

\vspace{2pt}
\textit{We name these four stages formally next class. For now: your objection to
the poster was not about arithmetic --- it was about where the number came from.
That is what this whole unit is for.}
\end{notesbox}

\boxguard
\begin{practicebox}
\small
Use the Student Life Survey table from section 2.

\textbf{(a)} Find the mean commute time for the six students.

\begin{work}
  \bar{x} &= \dfrac{20+8+35+15+5+25}{6} \\
          &= \dfrac{108}{6} \\
          &= 18 \text{ minutes}
\end{work}

\textbf{(b)} Explain why finding a ``mean lunch choice'' would \emph{not} make
sense.
\writeline
\writeline

\textbf{(c)} Lunch choice is categorical, so we summarize it with a percent
instead. What percent of these six students eat in the cafeteria?

\begin{work}
  \dfrac{3}{6} &= 0.5 \\
               &= 50\%
\end{work}

\textbf{(d)} Write one sentence interpreting your answer to (a). Name the group
and the units.
\writeline
\end{practicebox}

\end{document}
